
\begin{center}
{\huge \textbf{Citadelle marine}}\\ \vspace{0.5em}

\renewcommand{\arraystretch}{1.40}
\begin{tabular}{|C{5em}|C{40em}|}
\hline
Cadre  & Le plan d'eau est séparé en deux parties par une ligne matérialisée soit par des amers soit par des bouées de couleurs lestées. Les équipes doivent rester à une distance visible les unes des autres. \\
\hline
Règles  & Chaque équipe occupe une des deux parties du terrain. Elle localise le \textit{chmorblouf} situé sur son terrain (le \textit{chmorblouf} est matérialisé par un seau relié à un pare battage choisi pour sa couleur vive). Pour gagner la partie il faut réunir les deux \textit{chmorbloufs} dans son terrain. Une équipe peut déplacer son propre \textit{chmorblouf} uniquement quand aucun bateau ennemi est dans sa zone et à condition de ne pas l'ammener à plus d'un demi Mille de la ligne.
\\
\hline
Equipes & Deux bateaux par équipes, les équipes sont distingués grâce aux pavillons de couleur hissé dans la balancine. Un ou une chef.fe est placé.e sur chaque bateau. 
\\
\hline
Imaginaire & Les explorateurs du ciel zoologistes sont à la recherche des \textit{chmorbloufs}, des êtres légendaires vivants au creux des nuages dont le rôle écosystémique est majeur mais qui sont malheureusement en voie de disparition. Les \textit{chmorbloufs} ont un problème, leur sens de l'orientation est très mauvais et lors de la saison des amours ils peinent à se retrouver. Le but des zoologistes est de les aider dans cette tâche. Le prestige de l'équipe de zoologistes qui réussira la première cet exploit sera immense ... \\
\hline
Rôles & Le PE présent sur chaque bateau veille à la sécurité du bateau et de l'équipage, il reste pleinement vigilant à son environnement. 
Le ou la chef.fe présent.e sur chaque bateau s'assure du respect des règles de sécurité. 
Le CF garde un vision d'ensemble de la flotille, s'il lui semble qu'un bateau ne respecte pas les règles de sécurité il n'hésite pas à interrompre le jeu. \\
\hline
Action & Un bateau présent dans la zone adverse peut se faire éliminer s'il est atteint en sa voile ou en sa coque par une pomme de tooling adverse (constituée avec une garcette de diamètre d'environ un cm et d'une longueur d'environ 20m). 
Un bateau touché doit rentrer dans sa zone et libérer son \textit{chmorblouf} s'il en transportait un. 
 \\
\hline
Sens & Prendre en main de manière fine l'évolution du bateau sur le plan d'eau. 
Apprendre à régler finement ses voiles. 
Savoir se repérer sur un plan d'eau et contrôler sa trajectoire pour atteindre un objectif. Atteindre un point fixe sur l'eau à une vitesse modérée. 
Sensibiliser à la protection des écosystèmes. \\
\hline
Sécurité & La quête du \textit{chmorblouf} est impitoyable mais elle doit s'opérer en toute sécurité. 
Elle ne peut se faire que par temps calme sur une mer peu formée. 
Le CF doit avoir une pleine confiance en ses PE et dans les capacités des équipages en général. 
Un rappel des règles de priorité à la voile sera fait à l'avance.
On convient d'un signal d'arrêt et d'un point de rendez-vous si les équipes se perdent.  
Les \textit{chmorbloufs} sont des êtres sensibles qu'il faut manipuler avec précaution, sur chaque bateau un seul équipier est formé à cette tâche et sera responsable de la gaffe. 
Cette équipier doit porter une longe attachée au bateau en permanence, si l'équipier à la gaffe change la longe aussi. 
Le ou la cheffe présent.e sur chaque bateau s'assure du bon respect de cette règle.   
\\
\hline
\end{tabular}
\end{center}

\vspace{1cm}
